\documentclass[12pt]{article}
\usepackage[utf8]{inputenc}
\usepackage[T1]{fontenc}
\usepackage{hyperref}
\title{Kaburi Language Resources for K-12 Curriculum Development}
\author{Compiled by Claude Sonnet 5 (claude-sonnet-5) for Andr\'as Kornai}
\date{}
\begin{document}
\maketitle
\section{Introduction}
Kaburi (also Awe; ISO 639-3 uka) is a South Bird's Head language of Southwest Papua, Indonesia. Wikipedia gives 600 speakers (1986), while Joshua Project puts the Kaburi people at 1,100. Its classification is debated: some sources place it in Trans-New Guinea and others do not.
\section{Primary Educational Resources}
No primers, textbooks or school programmes were found in the search results.
\section{Linguistic and Reference Materials}
\begin{itemize}
\item Wikipedia, Kaburi language: \url{https://en.wikipedia.org/wiki/Kaburi_language}
\item Wikipedia, South Bird's Head languages: \url{https://en.wikipedia.org/wiki/South_Bird%27s_Head_languages}
\item Leiden University scholarly publication surfaced in the search (its content was not verified): \url{https://scholarlypublications.universiteitleiden.nl/access/item%3A3463909/view}
\item ASJP word list entry: \url{https://asjp.clld.org/languages/KABURI}
\end{itemize}
\section{Religious and Community Resources}
\begin{itemize}
\item Joshua Project profile of Kaburi in Indonesia: \url{https://www.joshuaproject.net/people_groups/12397/ID}
\end{itemize}
\section{Digital and Audio Resources}
No audio or video resources were found.
\section{Curriculum Development Context}
Documentation is thin and no teaching materials were found. Curriculum work would start from the academic sources above and from community input.
\end{document}